% Options for packages loaded elsewhere
\PassOptionsToPackage{unicode}{hyperref}
\PassOptionsToPackage{hyphens}{url}
\documentclass[twocolumn]{article}
\usepackage[margin=0.1in,top=0.1in,bottom=0.1in]{geometry}
\usepackage{xcolor}
\usepackage{amsmath,amssymb}
\usepackage{ragged2e}
\setcounter{secnumdepth}{-\maxdimen} % remove section numbering
\usepackage{iftex}
\ifPDFTeX
  \usepackage[T1]{fontenc}
  \usepackage[utf8]{inputenc}
  \usepackage{textcomp} % provide euro and other symbols
\else % if luatex or xetex
  \usepackage{unicode-math} % this also loads fontspec
  \defaultfontfeatures{Scale=MatchLowercase}
  \defaultfontfeatures[\rmfamily]{Ligatures=TeX,Scale=1}
\fi
\usepackage{lmodern}
\ifPDFTeX\else
  % xetex/luatex font selection
\fi
% Use upquote if available, for straight quotes in verbatim environments
\IfFileExists{upquote.sty}{\usepackage{upquote}}{}
\IfFileExists{microtype.sty}{% use microtype if available
  \usepackage[]{microtype}
  \UseMicrotypeSet[protrusion]{basicmath} % disable protrusion for tt fonts
}{}
\makeatletter
\@ifundefined{KOMAClassName}{% if non-KOMA class
  \IfFileExists{parskip.sty}{%
    \usepackage{parskip}
  }{% else
    \setlength{\parindent}{0pt}
    \setlength{\parskip}{6pt plus 2pt minus 1pt}}
}{% if KOMA class
  \KOMAoptions{parskip=half}}
\makeatother
\setlength{\emergencystretch}{3em} % prevent overfull lines
\providecommand{\tightlist}{%
  \setlength{\itemsep}{\GlobalItemSep}\setlength{\parskip}{\GlobalListParSkip}}
\usepackage{bookmark}
\IfFileExists{xurl.sty}{\usepackage{xurl}}{} % add URL line breaks if available
\urlstyle{same}
\hypersetup{
  hidelinks,
  pdfcreator={LaTeX via pandoc}}

% --------- Compression knobs (tune these values) ---------
\newcommand{\BodyFontSize}{\fontsize{5.2}{6.2}\selectfont}
\newcommand{\SectionFontStyle}{\fontsize{6.4}{7.2}\selectfont\bfseries}
\newif\ifUseRaggedRight
\UseRaggedRighttrue
\newif\ifUseMathabxSqbullet
\UseMathabxSqbullettrue
\newlength{\SectionBeforeSkip}
\newlength{\SectionAfterSkip}
\newlength{\GlobalItemSep}
\newlength{\GlobalListParSkip}
\newlength{\GlobalListTopSep}
\newlength{\ListLabelSep}
\newlength{\ListIndentLevelOne}
\newlength{\ListIndentLevelTwo}
\newlength{\ListIndentLevelThreePlus}
\setlength{\SectionBeforeSkip}{0.35ex plus 0.1ex minus 0.05ex}
\setlength{\SectionAfterSkip}{0.2ex plus 0.05ex}
\setlength{\GlobalItemSep}{0pt}
\setlength{\GlobalListParSkip}{0pt}
\setlength{\GlobalListTopSep}{0pt}
\setlength{\ListLabelSep}{0.15em}
\setlength{\ListIndentLevelOne}{0.5em}
\setlength{\ListIndentLevelTwo}{0.5em}
\setlength{\ListIndentLevelThreePlus}{0.5em}
\setlength{\columnsep}{8pt}

% Deep nested bullet symbol (level 3+): prefer \sqbullet if available.
\newcommand{\DeepLevelBullet}{$\scriptscriptstyle\blacksquare$}
\ifUseMathabxSqbullet
  \IfFileExists{mathabx.sty}{
    \usepackage{mathabx}
    \renewcommand{\DeepLevelBullet}{$\sqbullet$}
  }{}
\fi

% Math spacing knobs (horizontal spacing in equations)
% Tighten these further if gaps around '=' and '+' still look too large.
\thinmuskip=1mu
\medmuskip=1.5mu plus 0.5mu minus 0.5mu
\thickmuskip=2mu plus 0.5mu minus 0.5mu

\makeatletter
\renewcommand\section{\@startsection{section}{1}{\z@}{\SectionBeforeSkip}{\SectionAfterSkip}{\normalfont\SectionFontStyle}}
\makeatother

% Underfull \hbox fix for narrow columns: avoids stretched interword spaces.
\ifUseRaggedRight
  \AtBeginDocument{\RaggedRight}
\fi

% Compatibility: source notes use \bold, \argmax, and \argmin in math mode.
\providecommand{\bold}[1]{\mathbf{#1}}
\DeclareMathOperator*{\argmax}{arg\,max}
\DeclareMathOperator*{\argmin}{arg\,min}

% Pandoc conversion creates deeply nested lists; extend itemize depth.
\IfFileExists{enumitem.sty}{
  \usepackage{enumitem}
  \setlistdepth{9}
  \renewlist{itemize}{itemize}{9}
  \setlist[itemize]{itemsep=\GlobalItemSep,parsep=\GlobalListParSkip,topsep=\GlobalListTopSep,partopsep=0pt,labelsep=\ListLabelSep}
  \setlist[itemize,1]{label=\textbullet,leftmargin=\ListIndentLevelOne}
  \setlist[itemize,2]{label=$\circ$,leftmargin=\ListIndentLevelTwo}
  \setlist[itemize,3]{label=\DeepLevelBullet,leftmargin=\ListIndentLevelThreePlus}
  \setlist[itemize,4]{label=\DeepLevelBullet,leftmargin=\ListIndentLevelThreePlus}
  \setlist[itemize,5]{label=\DeepLevelBullet,leftmargin=\ListIndentLevelThreePlus}
  \setlist[itemize,6]{label=\DeepLevelBullet,leftmargin=\ListIndentLevelThreePlus}
  \setlist[itemize,7]{label=\DeepLevelBullet,leftmargin=\ListIndentLevelThreePlus}
  \setlist[itemize,8]{label=\DeepLevelBullet,leftmargin=\ListIndentLevelThreePlus}
  \setlist[itemize,9]{label=\DeepLevelBullet,leftmargin=\ListIndentLevelThreePlus}
}{
  % If enumitem is not installed, LaTeX may fail on very deep nesting.
}

\author{}
\date{}

\begin{document}
\BodyFontSize

\section{Markov Decision Processes and Policy / Value
Iteration}\label{markov-decision-processes-and-policy--value-iteration}

\begin{itemize}
\tightlist
\item
  \textbf{Markov Decision Processes} model environments where
  \emph{future states} are independent of \emph{past states} given the
  \emph{present state}

  \begin{itemize}
  \tightlist
  \item
    \(P^\pi(s' | s) = \sum_{a \in A} \pi(a|s)P(s'|s, a)\)
  \item
    \(R^\pi(s) = \sum_{a \in A}\pi(a|s)R(s, a)\)
  \end{itemize}
\item
  \textbf{Value Function} is the expected return starting from state
  \(s\) and following policy \(\pi\) while the \textbf{Action-Value
  Function} \(q^\pi(s,a)\) is the expected return starting from state
  \(s\) and taking action \(a\), and then following policy \(\pi\)

  \begin{itemize}
  \tightlist
  \item
    \(v^\pi(s) = \sum_{a \in A}\pi(a|s)q^\pi(s, a)\)
  \item
    \(q^\pi(s, a) = R^a_s + \gamma \sum_{s' \in S}P(s' |s ,a) v^\pi(s')\)
  \item
    Bellman Relationships:

    \begin{itemize}
    \tightlist
    \item
      \(v^\pi(s) = \mathbb{E}_\pi[R_{t+1} + \gamma v^\pi (s_{t+1}) | s_t = s] = \sum_{a \in A}\pi(a|s)(R(s, a) + \gamma \sum_{s' \in S} P(s' | s, a) v^\pi (s'))\)

      \begin{itemize}
      \tightlist
      \item
        Intuition: Marginalize over all possible actions (based on
        distribution from policy) - getting the expected return based on
        the immediate reward as well as the discounted future reward
        (which we must base on all possible next states given current
        state and the marginalized action)
      \end{itemize}
    \item
      \(q^\pi(s,a) = \mathbb{E}_\pi[R_{t+1} + \gamma q^\pi (s_{t+1}, A_{t+1}) | s_t = s, A_t = a] = R(s, a) + \gamma \sum_{s' \in S}P(s' | s, a) \sum_{a' \in A} \pi(a' | s')q^\pi (s', a')\)

      \begin{itemize}
      \tightlist
      \item
        Intuition: The immediate action is given (as a parameter), so we
        just need to marginalize the \emph{next state} given the current
        state and chosen action; from there, we average over all
        action-value functions as well
      \end{itemize}
    \end{itemize}
  \end{itemize}
\item
  Given a policy and the full environment dynamics, the value function
  associated with that policy can be computed via dynamic programming -
  \textbf{synchronous backup}

  \begin{itemize}
  \tightlist
  \item
    \(v_{t+1}(s) = \sum_{a \in A}\pi(a|s) (R(s, a) + \gamma \sum_{s' \in S} P(s' | s, a) v_t (s'))\)
  \end{itemize}
\item
  To compute the \textbf{optimal policy}, we can maximize over the
  optimal state-action value function

  \begin{itemize}
  \tightlist
  \item
    \textbf{Policy Iteration}: Iterate between evaluating the policy
    \(\pi\) and improving it greedily with respect to \(v^\pi\)

    \begin{itemize}
    \tightlist
    \item
      Evaluation: Use synchronous backup to find \(v^\pi\)
    \item
      Improvement:

      \begin{itemize}
      \tightlist
      \item
        Compute
        \(q^{\pi_i}(s, a) = R(s, a) + \gamma \sum_{s' \in S}P(s' | s, a) v^{\pi_i}(s')\)

        \begin{itemize}
        \tightlist
        \item
          This gives a \emph{q-table}, where there is an \emph{expected
          reward} for each action in each state
        \end{itemize}
      \item
        Compute \(\pi_{i + 1}(s) = \argmax_{a} q^{\pi_i}(s, a)\)
      \end{itemize}
    \item
      Repeat until convergence (policy does not change) - this means
      \(q^\pi(s, \pi(s)) = v^\pi(s)\)
    \end{itemize}
  \item
    \textbf{Value Iteration}: Iterate directly to find the optimal
    \(v^\pi\), and then extract the policy once afterwards

    \begin{itemize}
    \tightlist
    \item
      Repeat:

      \begin{itemize}
      \tightlist
      \item
        \(q_{k + 1}(s, a) = R(s, a) + \gamma \sum_{s' \in S}P(s' | s, a) v_k(s')\)
      \item
        \(v_{k + 1}(s) = \max_{a} q_{k+1}(s,a)\)
      \end{itemize}
    \item
      After convergence:

      \begin{itemize}
      \tightlist
      \item
        \(\pi(s) = \argmax_{a} [R(s, a) + \gamma \sum_{s' \in S}P(s' | s, a) v_{k + 1}(s')]\)
      \end{itemize}
    \end{itemize}
  \end{itemize}
\item
  The aforementioned algorithms \emph{only work if the model
  (environment) dynamics are known}
\end{itemize}

\section{Model-Free Prediction and
Control}\label{model-free-prediction-and-control}

\begin{itemize}
\tightlist
\item
  \textbf{Monte Carlo Policy Evaluation}: Sample trajectories, compute
  the returns, and average them to estimate \(v^\pi(s)\)

  \begin{itemize}
  \tightlist
  \item
    This works well, even in known environments, if transition
    probabilities are too complex to compute
  \item
    Algorithm:

    \begin{itemize}
    \tightlist
    \item
      Loop:

      \begin{itemize}
      \tightlist
      \item
        Generate an episode following \(\pi\):
        \(S_0, A_0, R_1, S_1, A_1, R_2, ..., S_{T_1}, A_{T_1}, R_T\)
      \item
        Loop for each step of episode backwards:
        \(t = T - 1, T - 2, ... 0\):

        \begin{itemize}
        \tightlist
        \item
          \(G \leftarrow \gamma G + R_{t+1}\)
        \item
          Append \(G\) to \(Returns(S_t)\)
        \item
          \(V(S_t) \leftarrow \text{average}(Returns(S_t))\)
        \end{itemize}
      \end{itemize}
    \end{itemize}
  \item
    Can use an incremental mean:
    \(\mu_t = \frac{1}{t} \sum_{j=1}^t x_j = \mu_{t-1} + \frac{1}{t}(x_t - \mu_{t-1})\)
  \end{itemize}
\item
  \textbf{Temporal Difference Learning}: Uses \emph{boostrapping}, which
  involves updating an estimate using \emph{other estimates} (like in
  dynamic programming) rather than waiting for the whole episode to
  complete. This is advantageous in that it can learn from
  \emph{incomplete} sequences and \emph{non-terminating} environments

  \begin{itemize}
  \tightlist
  \item
    \textbf{TD(0)}:
    \(v(S_t) \leftarrow v(S_t) + \alpha (R_{t+1} + \gamma v(S_{t+1}) - v(S_t))\)

    \begin{itemize}
    \tightlist
    \item
      View \([R_{t+1} + \gamma v(S_{t+1})] - v(S_t)\) as the
      \emph{error} of the existing estimate, which acts as a correction
    \end{itemize}
  \item
    \textbf{n-step}:
    \(G_t^n = R_{t+1} + \gamma R_{t+2} + ... + \gamma^{n-1}R_{t+n} + \gamma^n v(S_{t+n})\)

    \begin{itemize}
    \tightlist
    \item
      \(v(S_t) \leftarrow v(S_t) + \alpha(G^n_t - v(S_t))\)
    \end{itemize}
  \end{itemize}
\item
  \textbf{Generalized Policy Iteration}:

  \begin{itemize}
  \tightlist
  \item
    \textbf{Monte Carlo with Exploring Starts}:

    \begin{itemize}
    \tightlist
    \item
      Loop Forever:

      \begin{itemize}
      \tightlist
      \item
        Choose \(S_0 \in \mathcal{S}, A_0 \in \mathcal{A}(S_0)\)
        randomly such that all pairs have probability \textgreater{} 0
      \item
        Generate an episode from \(S_0, A_0\) following \(\pi\):
        \(S_0, A_0, R_1, ..., S_{T-1}, A_{T-1}, R_t\)
      \item
        Loop for each step of episode backwards:
        \(t = T - 1, T - 2, ..., 0\):

        \begin{itemize}
        \tightlist
        \item
          \(G \leftarrow \gamma G + R_{t+1}\)
        \item
          Append \(G\) to \(Returns(S_t, A_t)\)
        \item
          \(Q(S_t, A_t) \leftarrow \text{average}(Returns(S_t, A_t))\)
        \item
          \(\pi(S_t) \leftarrow \argmax_{a} Q(S_t, a)\)
        \end{itemize}
      \end{itemize}
    \end{itemize}
  \item
    \textbf{Monte Carlo with Epsilon-Greedy Exploration}:

    \begin{itemize}
    \tightlist
    \item
      \(\epsilon\)-Greedy Exploration puts probability \(1 - \epsilon\)
      to choose the greedy (optimal) action and probability \(\epsilon\)
      to choose a random action - the goal is to ensure there is
      sufficient exploration being done rather than exploiting a
      (potentially bad) policy
    \item
      Algorithm:

      \begin{itemize}
      \tightlist
      \item
        Loop:

        \begin{itemize}
        \tightlist
        \item
          Sample \((S_1, A_1, R_2, ..., S_T) \sim \pi_k\)
        \item
          For each \(S_t\) and \(A_t\):

          \begin{itemize}
          \tightlist
          \item
            \(N(S_t, A_t) \leftarrow N(S_t, A_t) + 1\)
          \item
            \(Q(S_t, A_t) \leftarrow Q(S_t, A_t) + \frac{1}{N(S_t, A_t)}(G_t - Q(S_t, A_t))\)

            \begin{itemize}
            \tightlist
            \item
              This is the incremental mean
            \end{itemize}
          \end{itemize}
        \item
          \(k \leftarrow k + 1\), \(\epsilon \leftarrow 1 / k\)

          \begin{itemize}
          \tightlist
          \item
            Anneal \(\epsilon\)
          \end{itemize}
        \item
          \(\pi_k = \epsilon-\text{greedy}(Q)\)
        \end{itemize}
      \end{itemize}
    \end{itemize}
  \item
    \textbf{On-Policy TD Control (SARSA)}: Estimates \(Q\) using TD

    \begin{itemize}
    \tightlist
    \item
      Repeat for each episode:

      \begin{itemize}
      \tightlist
      \item
        Initialize \(S\)
      \item
        Choose \(A\) from \(S\) using \(\epsilon\)-greedy policy
      \item
        Repeat (for each step of episode):

        \begin{itemize}
        \tightlist
        \item
          Take action \(A\), observe \(R\), \(S'\)
        \item
          Choose \(A'\) from \(S'\) using \(\epsilon\)-greedy
        \item
          \(Q(S, A) \leftarrow Q(S,A) + \alpha[R + \gamma Q(S', A') - Q(S, A)]\)

          \begin{itemize}
          \tightlist
          \item
            This can be generalized to n-step as well:
            \(q_t^{(n)}= R_{t+1} + \gamma R_{t+2} + ... + \gamma^{n-1}R_{t+n} + \gamma^n Q(S_{t+n}, A_{t+n})\)
          \end{itemize}
        \item
          \(S \leftarrow S'; A \leftarrow A'\)
        \end{itemize}
      \end{itemize}
    \end{itemize}
  \item
    \textbf{Off-Policy TD Control (Q-Learning)}: Use \emph{two different
    policies} - one being learned about and becomes optimal policy, and
    another which is more exploratory (used to generate trajectories);
    doing this allows for reuse of old experiences

    \begin{itemize}
    \tightlist
    \item
      Repeat for each episode:

      \begin{itemize}
      \tightlist
      \item
        Initialize \(S\)
      \item
        Repeat (for each step of episode):

        \begin{itemize}
        \tightlist
        \item
          Choose \(A\) from \(S\) using \(\epsilon\)-greedy derived from
          \(Q\)

          \begin{itemize}
          \tightlist
          \item
            This is the \textbf{behavior policy}, which is
            epsilon-greedy on the target \(Q\)
          \end{itemize}
        \item
          Take action \(A\), observe \(R\), \(S'\)
        \item
          \(Q(S, A) \leftarrow Q(S,A) + \alpha[R + \gamma \max_a Q(S', a) - Q(S, A)]\)

          \begin{itemize}
          \tightlist
          \item
            This is the \textbf{target policy}, which is fully greedy on
            \(Q\)
          \item
            The learning \textquotesingle acts\textquotesingle{} as if
            it took the target policy when updating, but the steps
            actually taken in the environment solely use the behavior
            policy
          \end{itemize}
        \item
          \(S \leftarrow S'\)
        \end{itemize}
      \end{itemize}
    \end{itemize}
  \end{itemize}
\item
  \textbf{Importance Sampling}: Rather than sampling from a given
  distribution, it may be easier to sample from a \emph{different
  distribution} and then correct the weight (e.g. easier to sample from
  behavior policy than target)

  \begin{itemize}
  \tightlist
  \item
    \(\mathbb{E}_{\tau \sim \pi}[r(\tau)] = \int P_\pi(\tau) r(\tau) d\tau = \int P_\mu (\tau) \frac{P_\pi (\tau)}{P_\mu (\tau)} r(\tau) d\tau = \mathbb{E}_{\tau \sim \mu}[\frac{P_\pi (\tau)}{P_\mu (\tau)} r(\tau)] \approx \frac{1}{n} \sum_i \frac{P_\pi (\tau_i)}{P_\mu (\tau_i)} r(\tau_i)\)
  \item
    This can be used to perform off-policy Monte Carlo and TD Learning;
    use the behavior policy to generate an episode or a TD target, and
    then reweight the return or TD target to use for the update

    \begin{itemize}
    \tightlist
    \item
      e.g. TD:

      \begin{itemize}
      \tightlist
      \item
        \(V(S_t) \leftarrow V(S_t) + \alpha(\frac{\pi(A_t | S_t)}{\mu(A_t | S_t)}(R_{t+1} + \gamma V(S_{t+1})) - V(S_t))\)
      \end{itemize}
    \end{itemize}
  \end{itemize}
\end{itemize}

\section{Value Function Approximation and Deep
Q-Learning}\label{value-function-approximation-and-deep-q-learning}

\begin{itemize}
\tightlist
\item
  Same approaches as model-free prediction work with function
  approximators

  \begin{itemize}
  \tightlist
  \item
    Monte Carlo Value Function Approximation:
    \(\Delta \bold{w} = \alpha (G_t - \hat{v}(S_t, \bold{w})) \nabla_{\bold{w}} \hat{v}(S_t, \bold{w})\)

    \begin{itemize}
    \tightlist
    \item
      \(G_t\) is \emph{unbiased}, but \emph{noisy} sample of true value
    \end{itemize}
  \item
    Temporal Difference Value Function Approximation:
    \(\Delta \bold{w} = \alpha (R_{t+1} + \gamma \hat{v}(S_{t+1}, \bold{w}) - \hat{v}(S_t, \bold{w})) \nabla_{\bold{w}} \hat{v}(S_t, \bold{w})\)

    \begin{itemize}
    \tightlist
    \item
      \(R_{t+1} + \gamma \hat{v}(S_{t+1}, \bold{w})\) is \emph{biased}
      estimate of true value since it is drawn from previous estimate
    \item
      The gradient ascent update is an example of \textbf{semi-gradient}
      learning
    \end{itemize}
  \end{itemize}
\item
  The model-free control algorithms are \emph{also} similar -
  approximate \(q\)

  \begin{itemize}
  \tightlist
  \item
    MC:
    \(\Delta \bold{w} = \alpha (G_t - \hat{q}(S_t, A_t, \bold{w})) \nabla_{\bold{w}} \hat{q} (S_t, A_t, \bold{w})\)
  \item
    Sarsa:
    \(\Delta \bold{w} = \alpha (R_{t+1} + \gamma \hat{q}(S_{t+1}, A_{t+1}, \bold{w}) - \hat{q}(S_t, A_t, \bold{w})) \nabla_{\bold{w}} \hat{q} (S_t, A_t, \bold{w})\)
  \item
    Q-Learning:
    \(\Delta \bold{w} = \alpha (R_{t+1} + \gamma \max_a \hat{q}(S_{t+1}, a, \bold{w}) - \hat{q}(S_t, A_t, \bold{w})) \nabla_{\bold{w}} \hat{q} (S_t, A_t, \bold{w})\)
  \end{itemize}
\item
  Linear Least-Squares Prediction Algorithms:

  \begin{itemize}
  \tightlist
  \item
    LSMC:
    \(0 = \sum_{t=1}^T \alpha(G_t - \hat{v}(S_t, \bold{w}))\bold{x}(S_t)\)

    \begin{itemize}
    \tightlist
    \item
      \(\bold{w} = (\sum_{t=1}^T \bold{x}(S_t) \bold{x}(S_t)^T)^{-1} \sum_{t=1}^T \bold{x}(S_t)G_t\)
    \end{itemize}
  \item
    LSTD:
    \(0 = \sum_{t=1}^T \alpha(R_{t+1} + \gamma \hat{v}(S_{t+1}, \bold{w}) - \hat{v}(S_t, \bold{w}))\bold{x}(S_t)\)

    \begin{itemize}
    \tightlist
    \item
      \(\bold{w} = (\sum_{t=1}^T \bold{x}(S_t) \bold{x}(S_t)^T - \gamma \bold{x}(S_{t+1}))^{-1} \sum_{t=1}^T \bold{x}(S_t)R_{t+1}\)
    \end{itemize}
  \item
    LSDQ:
    \(0 = \sum_{t=1}^T \alpha(R_{t+1} + \gamma \hat{q} (S_{t+1}, \pi (S_{t+1}), \bold{w}) - \hat{q} (S_t, A_t, \bold{w}))\bold{x} (S_t, A_t)\)

    \begin{itemize}
    \tightlist
    \item
      \(\bold{w} = (\sum_{t=1}^T \bold{x} (S_t, A_t)(\bold{x} (S_t, A_t) - \gamma \bold{x} (S_{t+1}, \pi (S_{t+1})))^T)^{-1} \sum_{t=1}^T \bold{x} (S_t, A_t) R_{t+1}\)
    \item
      Sample \(S_t, A_t, R_{t+1}, S_{t+1}\) from old policy to generate
      an experience replay
    \end{itemize}
  \end{itemize}
\item
  Deep Q-Learning: Standard Q-Learning approach, \emph{however},
  incorporates (1) \textbf{Experience replay} to address correlation
  between samples and (2) \textbf{Fixed Q targets} to handle
  non-stationary targets

  \begin{itemize}
  \tightlist
  \item
    Experience Replay: Store \((s_t, a_t, r_t, s_{t+1})\) suples in
    replay memory, and then sample an experience tuple (or batch) from
    this memory to compute the target value and prediction used to
    perform gradient descent
  \item
    Fixed Targets: Use \(\bold{w^-}\) as the set of weights used in the
    \emph{target} and \(\bold{w}\) as the set of weights being updated.
    Periodically adjust \(\bold{w^-} \leftarrow \bold{w}\) (frequency is
    hyperparameter)

    \begin{itemize}
    \tightlist
    \item
      \(\Delta \bold{w} = \alpha (r + \gamma \max_a \hat{Q}(s', a', \bold{w^-}) - \hat{Q}(s, a, \bold{w}))\nabla_{\bold{w}} \hat{Q}(s, a, \bold{w})\)
    \end{itemize}
  \item
    Some improvements: \textbf{Double DQN}: use \emph{two networks},
    where one performs the action selection and another creates the
    target value (\(\argmax \hat{Q}\) term); \textbf{Dueling DQN}:
    branch the same network where one branch generates \(V(s)\) and the
    other generates \(A(s, a)\) so that (\(Q(s, a) = A(s, a) + V(s)\)) -
    the goal is to encourage learning which states are valuable without
    needing to learn the effect of action at each state;
    \textbf{Prioritized Experience Replay}: Prioritize experiences where
    there is a big difference between prediction and TD target (can use
    a priority queue)
  \end{itemize}
\item
  Some issues that can lead to instability in training -
  \textbf{function approximators}, \textbf{bootstrapping}, and
  \textbf{off-policy training}
\item
  Batch-based methods are often used for training - this is more
  sample-efficient
\end{itemize}

\section{Policy Optimization I: Policy Gradient and
Actor-Critic}\label{policy-optimization-i-policy-gradient-and-actor-critic}

\begin{itemize}
\tightlist
\item
  Parameterize and optimize policy directly

  \begin{itemize}
  \tightlist
  \item
    Discrete Policy (\textbf{Softmax}):
    \(\pi_\theta (a | s) = \frac{\exp (\phi(s, a)^{T} \theta)}{\sum_{a'} \exp (\phi(s, a')^T \theta)}\)

    \begin{itemize}
    \tightlist
    \item
      \(\nabla_\theta \log \pi_\theta (a|s) = \phi(s,a) - \mathbb{E}_{\pi_\theta}[\phi(s, .)]\)
    \end{itemize}
  \item
    Continuous Policy (\textbf{Gaussian)}:
    \(\mu (s) = \phi(s)^T \theta\)

    \begin{itemize}
    \tightlist
    \item
      \(a \sim \mathcal{N} (\mu (s), \sigma^2)\)
    \item
      \(\nabla \theta \log(\pi_\theta) (a |s) = \frac{(a - \mu(s)) \phi(s)}{\sigma^2}\)
    \end{itemize}
  \end{itemize}
\item
  Gradient-Free Methods:

  \begin{itemize}
  \tightlist
  \item
    \textbf{Evolution Strategy}: Perturb \(\theta\) with multiple noises
    and execute rollouts to compute \(J(\theta)\) for each; set new
    \(\theta\) to weighted sum based on scores
    (\(\theta_{t+1} \leftarrow \theta_t + \alpha \frac{1}{n \delta} \sum_{i=1}^n J(\theta_i) \epsilon_i\))
  \item
    \textbf{Cross-Entropy Method}: Sample \(\theta\) from distribution,
    execute rollouts, and then update \(\mu\) (distribution parameter)
    to maximize the likelihood of the top n\% of \(\theta\) generated
    based on score
    (\(\mu^{(i+1)} = \argmax_\mu \sum_{k \in \hat{C}} \log P_\mu (\theta^{(k)})\))
  \end{itemize}
\item
  Log Ratio Trick:
  \(\nabla_\theta \pi_\theta (a|s) = \pi_\theta (a|s) \frac{\nabla_\theta \pi_\theta (a|s)}{\pi_\theta (a|s)} = \pi_\theta (a|s) \nabla_\theta \log \pi_\theta (a|s)\)
\item
  The objective of policy gradient is to maximize the reward from a
  trajectory sampled from the policy:
  \(J(\theta) = \mathbb{E}_{\pi_\theta}[\sum_{t=0}^T R(s_t, a_t)] = \sum_\tau P(\tau ; \theta) R(\tau)\)
  where
  \(\tau = (s_0, a_0, r_1, ..., s_{T-1}, a_{T-1}, r_T, s_T) \sim (\pi_\theta, P(s_{t+1} | s_t, a_t))\)

  \begin{itemize}
  \tightlist
  \item
    To get this in a differentiable form:

    \begin{itemize}
    \tightlist
    \item
      \(\nabla_\theta J(\theta) = \nabla_\theta \sum_\tau P(\tau ; \theta) R(\tau) = \sum_\tau P(\tau; \theta) R(\tau) \nabla_\theta \log P(\tau; \theta)\)
      (\textbf{log ratio trick})
    \item
      \(\approx \frac{1}{m}\sum_{i=1}^m R(\tau_i) \nabla_\theta \log P(\tau_i; \theta)\)
      (\textbf{estimate of expectation over \(\tau\)})
    \item
      \(= \frac{1}{m}\sum_{i=1}^m R(\tau_i) \nabla_\theta \log [\mu(s_0) \prod_{t=0}^{T-1} \pi_\theta (a_t | s_t) p(s_{t+1} | s_t, a_t)]\)
      (\textbf{decomposition of \(P\)})
    \item
      \(= \frac{1}{m}\sum_{i=1}^m R(\tau_i) \nabla_\theta [\log \mu(s_0) + \sum_{t=0}^{T-1} \log \pi_\theta (a_t | s_t) + \log p(s_{t+1} | s_t, a_t)]\)
      (\textbf{can eliminate transition dynamics since they are not a
      function of \(\theta\)})
    \item
      \(\nabla_\theta J(\theta) \approx \frac{1}{m} \sum_{i=1}^m R(\tau_i) \sum_{t=0}^{T-1} \nabla_\theta \log \pi_\theta (a^i_t | s^i_t)\)

      \begin{itemize}
      \tightlist
      \item
        This is shifting the distribution (\(\pi\)) through parameter
        \(\theta\) in a way to ensure future samples achieve higher
        scores based on the reward function - this is akin to
        \textbf{weighted maximum likelihood}
      \end{itemize}
    \end{itemize}
  \end{itemize}
\item
  \textbf{Temporal Causality}: Policy should not affect rewards from the
  past, so only sum \emph{future rewards}

  \begin{itemize}
  \tightlist
  \item
    \(\nabla_\theta J(\theta) = \mathbb{E}_\tau [\sum_{t=0}^{T-1} \nabla_\theta \log \pi_\theta (a_t | s_t) \sum_{t' = t}^{T-1} r_t'] \approx \frac{1}{m} \sum_{i=1}^m \sum_{t = 0}^{T-1} G_t^{(i)} \cdot \nabla_\theta \log \pi_\theta (a^i_t | s^i _t)\)
  \end{itemize}
\item
  \textbf{Baselines}: Reduce variance by subtracting a baseline

  \begin{itemize}
  \tightlist
  \item
    \(\nabla_\theta J(\theta) = \mathbb{E}_\tau [\sum_{t=0}^{T-1} (G_t - b) \cdot \nabla_\theta \log \pi_\theta (a_t | s_t)]\)
    where \(b = \mathbb{E} [r_t + r_{t+1} + ... + r_{T-1}]\) (expected
    return)
  \end{itemize}
\item
  \textbf{REINFORCE}: On-policy algorithm

  \begin{itemize}
  \tightlist
  \item
    Generate episode following \(\pi(\theta)\)
  \item
    For each step of episode:

    \begin{itemize}
    \tightlist
    \item
      Let \(G\) be the return from step \(t\) onwards
    \item
      \(\theta \leftarrow \theta + \alpha \gamma^t G \nabla_\theta \log \pi (A_t | S_t, \theta)\)
    \end{itemize}
  \end{itemize}
\item
  \textbf{Actor-Critic}: Instead of using \(G_t\), which is sampled from
  Monte Carlo, use a critic to estimate the action-value function. This
  implies \emph{two sets of parameters} - one for the actor (policy
  function), and one for the critic (value function)

  \begin{itemize}
  \tightlist
  \item
    \(\nabla_\theta J(\theta) = \mathbb{E}_{\pi_\theta} [\sum_{t=0}^{T-1} Q_w (s_t, a_t) \cdot \nabla_\theta \log \pi_\theta (a_t | s_t)]\)
  \item
    Critic parameters can be updated using any prior approach (Monte
    Carlo, TD, etc.)
  \item
    Baseline can be applied here as well - the state-value function can
    be used to form the baseline (\textbf{advantage function})
    \(A^{\pi, \gamma}(s,a) = Q^{\pi, \gamma}(s, a) - V^{\pi, \gamma}(s)\)

    \begin{itemize}
    \tightlist
    \item
      Rather than maintain \emph{two sets of parameters} for just the
      critic, it can be recognized that the TD error is an unbiased
      estimate of the advantage function:
    \item
      \(\mathbb{E}_{\pi_\theta}[\delta^{\pi_\theta} | s, a] = \mathbb{E}_{\pi_\theta} [r + \gamma V^{\pi_\theta}(s') |s, a] - V^{\pi_\theta}(s) = Q^{\pi_\theta}(s, a) - V^{\pi_\theta}(s) = A^{\pi_\theta}(s,a)\)
    \item
      \(\nabla_\theta J(\theta) = \mathbb{E}_{\pi_\theta}[\nabla_\theta \log \pi_\theta (a |s)\delta^{\pi_\theta}]\),
      where \(\delta_v = r + \gamma V_{\kappa}(s') - V_\kappa (s)\)

      \begin{itemize}
      \tightlist
      \item
        Just need to maintain weights for \(V\), which can be updated
        using MC, TD, etc. (e.g.
        \(\Delta \kappa = \alpha (G - V_\kappa (s)) \psi (s)\))
      \end{itemize}
    \end{itemize}
  \end{itemize}
\end{itemize}

\section{Policy Optimization II}\label{policy-optimization-ii}

\begin{itemize}
\tightlist
\item
  Turn Policy Gradient into an \emph{off-policy} algorithm via
  \textbf{importance sampling}, where the \emph{old policy} is sampled
  from (using replay buffer) rather than the current one; must
  incorporate a constraint to limit excessive changes between subsequent
  policies (creates a \textbf{trust region})

  \begin{itemize}
  \tightlist
  \item
    \(J(\theta) = \mathbb{E}_{s, a \sim \pi_\theta}[r(s, a)] = \mathbb{E}_{s, a \sim \hat{\pi}}[\frac{\pi_\theta(s,a)}{\hat{\pi}(s,a)}r(s,a)]\)
    subject to
    \(KL(\pi_{\theta_{\text{old}}} ||\pi_\theta) = \sum_a \pi_{\theta_{\text{old}}}(a|s) \log \frac{\pi_\theta (a|s)}{\pi_{\theta_{\text{old}}}(a|s)} \leq \delta\)
  \end{itemize}
\item
  \textbf{Natural Policy Gradient}: Use second-order optimization
  methods to solve aforementioned constrained objective

  \begin{itemize}
  \tightlist
  \item
    Collect trajectories \(\mathcal{D}_k\) on policy
    \(\pi_k = \pi(\theta_k)\)
  \item
    Estimate advantages \(\hat{A}_t^{\pi_k}\) using any method
  \item
    Estimate \(g = \nabla_\theta J_{\theta_t}(\theta)\) and
    \(H = \nabla^2_\theta KL(\theta || \theta_t)\), which can also be
    expressed as
    \(F=\mathbb{E}_{\pi_\theta(s,a)}[\nabla \log \pi_\theta (s, a) \nabla \log \pi_\theta (s, a)^T]\)
  \item
    Compute update:
    \(\theta_{k+1}=\theta_k + \sqrt{\frac{2\delta}{\hat{g}_k^T \hat{H}_k^{-1}\hat{g}_k}}\hat{H}_k^{-1}\hat{g}_k\)
  \item
    \textbf{Disadvantage}: \(H\) and \(H^{-1}\) can be very difficult to
    compute
  \end{itemize}
\item
  \textbf{Trust Region Policy Optimization}: Estimates \(x=H^{-1}g\) by
  solving \(Hx=g\), which is equivalent to solving
  \(\min_x \frac{1}{2}x^T H x - g^Tx\) - do this via conjugate gradient
  method

  \begin{itemize}
  \tightlist
  \item
    \textbf{Disadvantage}: Computing \(H\) is still expensive, and
    requires a large batch of rollouts to properly approximate \(H\)
  \end{itemize}
\item
  \textbf{Proximal Policy Optimization (PPO) with Adaptive KL Penalty}:
  Incorporate the KL constraint into the objective directly

  \begin{itemize}
  \tightlist
  \item
    Collect set of partial trajectories \(\mathcal{D}_k\) on
    \(\pi_k = \pi(\theta_k)\)
  \item
    Estimate advantages \(\hat{A}_t^{\pi_k}\) using any method
  \item
    \(\theta_{k+1} = \argmax_\theta \mathbb{E}_{\pi_{\theta_{\text{old}}}}[\frac{\pi_\theta (a_t | s_t)}{\pi_{\theta_{\text{old}}}(a_t | s_t)}A_t] - \beta \mathbb{E}_{\pi_{\theta_{\text{old}}}}[KL[\pi_{\theta_{\text{old}}}, \pi_\theta]]\)

    \begin{itemize}
    \tightlist
    \item
      Using SGD
    \end{itemize}
  \item
    if \(D_{KL}(\theta_{k+1}||\theta_k) \geq 1.5\delta\) then
    \(\beta_{k+1} = 2\beta_K\) else \(\beta_{k+1} = \beta_k / 2\)

    \begin{itemize}
    \tightlist
    \item
      If there are significant differences, then have a higher
      divergence penalty
    \end{itemize}
  \end{itemize}
\item
  \textbf{PPO with Clipping}: Axe using KL divergence and instead clip
  the advantage function directly to avoid too large of changes in
  probability ratio
  \(r_t(\theta) = \frac{\pi_\theta (a_t | s_t)}{\pi_{\theta_{\text{old}}}(a_t | s_t)}\)

  \begin{itemize}
  \tightlist
  \item
    \(L_t (\theta) = \min(r_t (\theta) \hat{A}_t, \text{clip}(r_t(\theta), 1 - \epsilon, 1 + \epsilon) \hat{A}_t)\)
  \end{itemize}
\item
  \textbf{Deep Deterministic Policy Gradient (DDPG)}: Extend DQN to
  \emph{continuous} action space. Still utilize replay buffer and target
  network logic as with DQN

  \begin{itemize}
  \tightlist
  \item
    \(a^* = \argmax_a Q^* (s, a) \approx \argmax_\theta Q_\phi (s, \mu_\theta(s))\)

    \begin{itemize}
    \tightlist
    \item
      Train policy \(\mu_\theta(s)\) to give action that maximizes
      \(Q_\phi(s,a)\)
    \end{itemize}
  \item
    \textbf{Q-Function Update}:
    \(\min \mathbb{E}_{s, r, s', d \sim \mathcal{D}}[(r + \gamma Q_{\phi_{\text{targ}}}(s', \mu_{\theta_{\text{targ}}}(s')) - Q_\phi (s,a))^2]\)
  \item
    \textbf{Policy Update}:
    \(\max_\theta \mathbb{E}_{s \sim \mathcal{D}}[Q_\phi (s, \mu_\theta(s))]\)
  \item
    \textbf{Disadvantage}: Tends to \emph{dramatically overestimate
    Q-values}
  \end{itemize}
\item
  \textbf{Twin Delayed DDPG (TD3)}: Addresses issues with DDPG

  \begin{itemize}
  \tightlist
  \item
    \textbf{Double-Q}: Learn two Q functions, \(Q_{\phi_1}\) and
    \(Q_{\phi_2}\) using the same aforementioned loss function, but have
    the target be the smaller of the two:
    \(y(r, s') = r + \gamma \min_{i=1,2} Q_{\phi_{i, targ}} (s', a_{TD3}(s'))\)

    \begin{itemize}
    \tightlist
    \item
      Note that there are still \emph{two target networks} for each Q
      function
    \end{itemize}
  \item
    \textbf{Target Policy Smoothing}:
    \(a_{TD3}(s') = \text{clip}(\mu_{\theta, \text{targ}}(s') + \text{clip}(\epsilon, -c, c), a_{low}, a_{high}), \epsilon \sim \mathcal{N}(0, \sigma)\)

    \begin{itemize}
    \tightlist
    \item
      This regularizes by avoiding incorrect sharp peak for some actions
      output by the policy
    \end{itemize}
  \end{itemize}
\item
  \textbf{Soft Actor-Critic (SAC)}: Incorporates \textbf{entropy
  regularization} to ensure there is a tradeoff between expected return
  and entropy in the policy - this encourages exploration

  \begin{itemize}
  \tightlist
  \item
    Q-Function Update:

    \begin{itemize}
    \tightlist
    \item
      \(L(\phi_i, \mathcal{D}) = \mathbb{E}[(Q_\phi(s,a) - y(r, s', d))^2]\)
    \item
      \(y(r, s') = r + \gamma(\min_{j=1,2} Q_{\phi_{targ, j}}(s', \hat{a}') - \alpha \log \pi_\theta (\hat{a}' |s'))\)
    \item
      \(\hat{a}' \sim \pi_\theta(\cdot | s')\)

      \begin{itemize}
      \tightlist
      \item
        No need to reparameterize here because the gradient is for
        \(\phi\), not \(\theta\)
      \end{itemize}
    \end{itemize}
  \item
    Policy Update:

    \begin{itemize}
    \tightlist
    \item
      \(\max_\theta E_{s \sim \mathcal{D}, \epsilon \sim \mathcal{N}} [\min_{j=1,2} Q_{\phi_j}(s, \hat{a}_\theta (s, \epsilon)) - \alpha \log \pi_\theta (\hat{a}_\theta (s, \epsilon) | s)]\)
    \item
      To avoid the dependency of sampling from the policy in the
      objective, we used a \textbf{reparameterization trick} to instead
      sample over noise

      \begin{itemize}
      \tightlist
      \item
        \(\hat{a}_\theta(s, \epsilon) = \tanh(\mu_\theta(s) + \sigma_\theta (s) \odot \epsilon), \epsilon \sim \mathcal{N}(0, I)\)
      \end{itemize}
    \end{itemize}
  \end{itemize}
\end{itemize}

\section{Model-Based Reinforcement
Learning}\label{model-based-reinforcement-learning}

\begin{itemize}
\tightlist
\item
  Learn a \textbf{model} that can be used to \emph{plan} a better policy
  - this is much more \emph{sample efficient}, but can be difficult to
  learn since now there are two sources of error

  \begin{itemize}
  \tightlist
  \item
    Model learns to predict \(S_{t+1}\) and \(R_{t+1}\) in a supervised
    fashion from collected trajectories
  \item
    Planning: Sample experiences from the model and then use model-free
    reinforcement learning algorithms (Monte Carlo, Sarsa, Q-Learning)
    to the sampled experiences to learn
  \end{itemize}
\item
  \textbf{Dyna-Q} (Model-Based Value Optimization):

  \begin{itemize}
  \tightlist
  \item
    Execute action \(A\) based on policy (e.g. \(\epsilon\)-greedy), and
    observe resultant \(R\) and \(S'\)
  \item
    \(Q(S,A) \leftarrow Q(S,A) + \alpha [R + \gamma \max_a Q(S', a) - Q(S, A)]\)
  \item
    Update model from collected \(R\) and \(S'\) (e.g. regression)
  \item
    Repeat \(n\) times:

    \begin{itemize}
    \tightlist
    \item
      \(S \leftarrow \text{random previously observed state}\);
      \(A \leftarrow \text{random action previously taken in S}\)
    \item
      \(R, S' \leftarrow \text{Model}(S,A)\)
    \item
      \(Q(S,A) \leftarrow Q(S,A) + \alpha [R + \gamma \max_a Q(S' ,a) - Q(S,A)]\)
    \end{itemize}
  \end{itemize}
\item
  \textbf{Model-Based Policy Optimization}:

  \begin{itemize}
  \tightlist
  \item
    \textbf{Linear-Quadratic Regulator} methods can be used to solve
    \(\min_{a_i, ..., a_t}\sum_{t=1}^T c(s_t, a_t)\) subject to
    \(s_t = f(s_{t-1}, a_{t-1})\), where \(c\) is the cost function (can
    use \emph{negative reward}) and \(s_t\) is the transition dynamics
  \item
    Algorithm (Planning Incorporated):

    \begin{itemize}
    \tightlist
    \item
      Run base policy \(\pi_0(a_t | s_t)\) to collect
      \(\mathcal{D} = \{(s, a, s')_i\}\)
    \item
      Loop every step:

      \begin{itemize}
      \tightlist
      \item
        Loop every N steps:

        \begin{itemize}
        \tightlist
        \item
          Learn dynamics model \(s' = p(s, a)\) to minimize
          \(\sum_i ||f(s_i, a_i) - s'_i||^2\)
        \end{itemize}
      \item
        Plan through \(f(s, a)\) to choose actions (use LQR method)
      \item
        Execute first planned action and observe resulting state
      \item
        Append \((s, a, s')\) to dataset \(\mathcal{D}\)
      \end{itemize}
    \end{itemize}
  \item
    Algorithm (Learning and Planning Incorporated):

    \begin{itemize}
    \tightlist
    \item
      Run base policy \(\pi_0(a_t | s_t)\) to collect
      \(\mathcal{D} = \{(s, a, s')_i\}\)
    \item
      Loop:

      \begin{itemize}
      \tightlist
      \item
        Learn dynamics model \(f(s, a)\) to minimize
        \(\sum_i ||f(s_i, a_i) - s'_i||^2\)
      \item
        Backpropagate through \(f(s, a)\) into policy to optimize
        \(\pi_\theta(a_t | s_t)\)
      \item
        Run \(\pi_\theta (a_t | s_t)\), appending visited \((s, a, s')\)
        to \(\mathcal{D}\)
      \end{itemize}
    \end{itemize}
  \end{itemize}
\end{itemize}

\section{Imitation Learning}\label{imitation-learning}

\begin{itemize}
\tightlist
\item
  \textbf{Imitation Learning} becomes a \emph{supervised learning} of
  the policy network, but standard behavioral cloning is limited since
  the policy can run into off-course situations
\item
  \textbf{Dataset Aggregation (DAgger)}: Initially collect human data
  and train \(\pi_\theta\) from human data. Then, run \(\pi_\theta\) to
  get a new dataset that is then \emph{labeled by humans with correct
  actions}. Finally, aggregate the newly labeled dataset with the
  original (already labeled dataset)

  \begin{itemize}
  \tightlist
  \item
    Human annotation is expensive - could try to substitute this out
    with an offline brute-force search (e.g. other, much slower
    algorithm)
  \end{itemize}
\item
  \textbf{HG-DAgger}: Same approach as dagger, but instead \emph{gate}
  human intervention - if human decides to take control, then record
  their state and actions into the dataset

  \begin{itemize}
  \tightlist
  \item
    One drawback here is that the bad agent actions are ignored in the
    original formulation:
    \(\mathcal{L}_{\text{HG-DAgger}(\pi)} = \mathbb{E}_{(s, a_h) \sim \mathcal{D}_{\text{takeover}}}[||\pi(s) - a_h||^2]\)
  \end{itemize}
\item
  \textbf{Proxy Value Propagation}: In cases of intervention, discourage
  agent action \(a_n\) while promoting human action \(a_h\)

  \begin{itemize}
  \tightlist
  \item
    \(J^{PV}(\theta) = \mathbb{E}_{(s, a_n, a_h)}[|Q_\theta (s, a_h) - 1|^2 + |Q_\theta(s, a_n) + 1|^2]\)
  \item
    \(J^{TD}(\theta) = \mathbb{E}_{(s, a, s')}[|Q_\theta(s,a) - \gamma \max_{a'} Q_\theta (s', a)|^2]\)
  \item
    \(J(\theta) = J^{PV}(\theta) + J^{TD}(\theta)\) (Learn policy via
    \(Q\) optimized from \(J\))
  \end{itemize}
\item
  \textbf{Predictive Preference Learning}: Provide a visualization of
  agent trajectory, and leverage preference learning on human
  trajectories when they decide to takeover

  \begin{itemize}
  \tightlist
  \item
    \(L(\pi) = \mathbb{E}_{s, a^+, a^-} [\log \sigma (\beta \log \frac{\pi(a^+ | s)}{\pi (a^- | s)})]\)
  \end{itemize}
\item
  \textbf{Adaptive Intervention Mechanism (AIM)}: Learn an uncertainty
  estimator such that the agent requests help when necessary (e.g. when
  \(Q_\theta (s, a_n) < \delta\) quantile of \(Q_\theta(s, a_n)\))
\item
  \textbf{Inverse Reinforcement Learning}: Provide expert data to learn
  a cost/reward function to explain expert behaviors

  \begin{itemize}
  \tightlist
  \item
    \textbf{Maximum Entropy IRL}: Observed trajectories:
    \(\mathcal{D} = \{\tau_i\}\), model reward linearly
    \(R(\tau) = \sum_t w^T \phi(s_t, a_t)\)

    \begin{itemize}
    \tightlist
    \item
      Optimize for distribution that matches expert behavior but is
      otherwise maximally uncertain:
      \(\max_{p(\tau)} H(p) = - \sum_\tau p(\tau) \log p(\tau)\) subject
      to \(\mathbb{E}_p [\phi(\tau)] = \hat{\phi}_{expert}\)
    \item
      Result: \(p(\tau | w) = \frac{1}{Z(w)}\exp(R(\tau))\) where
      \(Z(w) = \sum_\tau \exp(R(\tau))\) (expert trajectories
      expontentially more likely if rewards are high)
    \item
      \(\mathcal{L} (w) = \sum_{\tau \in \mathcal{D}}[R(\tau) - \log Z(w)]\)
    \end{itemize}
  \item
    \textbf{Generative Adversarial Imitation Learning (GAIL)}: Train a
    discriminator to predict whether a given state \(s\) and action
    \(s\) is sampled from the demonstrations \(M\) or generated by
    running the policy \(\pi\)

    \begin{itemize}
    \tightlist
    \item
      \(\argmin_D - \mathbb{E}_{d^{\mathcal{M}}(s, a)}[\log (D(s,a))] - \mathbb{E}_{d^\pi (s,a)}[\log (1 - D(s,a))]\)
    \item
      Train policy with \(r_t = -\log (1 - D(s_t, a_t))\)
    \item
      If there are state-only demonstrations, then
      \(\argmin_D - \mathbb{E}_{d^{\mathcal{M}}(s, s')}[\log (D(s, s'))] - \mathbb{E}_{d^\pi (s, s')}[\log (1 - D(s, s'))]\)
      and
      \(r(s_t, s_{t+1}) = \max [0, 1 - 0.25 (D(s_t, s_{t+1}) - 1)^2]\)
    \end{itemize}
  \end{itemize}
\item
  Common paradigm is to \textbf{use imitation learning for pretraining}
  and then \textbf{finetune with standard methods (e.g. policy
  gradient)}

  \begin{itemize}
  \tightlist
  \item
    To avoid forgetting demonstrations, try to leverage off-policy
    reinforcement learning so that the demonstrations can be injected
    into the off-policy samples (maybe with \emph{importance sampling})
  \end{itemize}
\end{itemize}

\section{System Design}\label{system-design}

\begin{itemize}
\tightlist
\item
  \textbf{Model Parallelism}: Different machines are responsible for
  computation in parts of a \emph{single network}
\item
  \textbf{Data Parallelism}: Different machines have the complete copy
  of the model, but get different portions of the data
\item
  Approaches to Parameter Updating in Distributed Systems:

  \begin{itemize}
  \tightlist
  \item
    \textbf{Parameter Averaging}: Distribute copy of parameters to each
    worker, then train each worker on subset of data, and set up the
    global parameters to average the parameters from each worker -
    repeat
  \item
    \textbf{Distributed Stochastic Gradient Descent}: Transfer just
    updates (gradient) to the parameter server, and server periodically
    transfers weights to worker machines
  \item
    \textbf{Synchronous Updates} have agents training in parallel but
    then wait until all updates are collected before enforcing on global
    network; \textbf{Asynchronous Updates} have the global network
    updated periodically by individual workers but not necessarily
    simulatenously
  \item
    \textbf{Hogwild}: Normally, lock-based SGD requires a lock on
    current state of parameters before updating; however, since updates
    are often sparse there is no need for the locks since there is still
    convergence
  \end{itemize}
\item
  \textbf{Distributed Architectures}:

  \begin{itemize}
  \tightlist
  \item
    Separate out replay memory, actors, parameter servers, and learners
  \end{itemize}
\end{itemize}

\end{document}
